\documentclass{article}
\usepackage{amsmath,amssymb,amsthm}
\usepackage{xcolor}
\newtheorem{theorem}{Theorem}
\newtheorem{lemma}{Lemma}
\newtheorem{assumption}{Assumption}
\newcommand{\E}{\mathbb{E}}
\newcommand{\norm}[1]{\left\|#1\right\|}
\newcommand{\grad}{\nabla}
\newcommand{\R}{\mathbb{R}}

\begin{document}

\section*{Algorithm Name}
\textbf{Accelerated Gradient Method with Classical Momentum Schedule}\\
(This is the Nesterov’s accelerated gradient (NAG) method where the momentum
parameter is chosen as $\displaystyle \beta_k=\frac{k-1}{k+2}$.)

%%%%%%%%%%%%%%%%%%%%%%%%%%%%%%%%%%%%%%%%%%%%%%%%%%%%%%%%%%%%%%%%%%%%%%%%%%%%
\section*{Mathematical Formulation}
Let $f:\R^d\to\R$ be the objective function.  
Initialize $x_{0}=y_{0}\in\R^{d}$ and set a stepsize $\alpha>0$.  
For $k=0,1,2,\dots$ perform
\[
\boxed{
\begin{aligned}
\beta_k &:= \frac{k-1}{k+2},\\[4pt]
y_k    &= x_k + \beta_k\,(x_k-x_{k-1}), \qquad (k\ge 1,\;x_{-1}:=x_0),\\[4pt]
x_{k+1}&= y_k - \alpha\,\grad f(y_k).
\end{aligned}}
\]
The sequence $\{x_k\}$ is the main iterate; $\{y_k\}$ is the
extrapolated point used to evaluate the gradient.

%%%%%%%%%%%%%%%%%%%%%%%%%%%%%%%%%%%%%%%%%%%%%%%%%%%%%%%%%%%%%%%%%%%%%%%%%%%%
\section*{Pseudocode}
\begin{verbatim}
Input:   stepsize α > 0, initial point x0 ∈ R^d, maxIter T
Set   x[-1] ← x0 ,  x[0] ← x0
for k = 0,…,T-1 do
    if k == 0 then
        β ← 0                 // because (k-1)/(k+2)= -1/2 but we set β0=0
    else
        β ← (k-1)/(k+2)
    end if
    y ← x[k] + β * (x[k] - x[k-1])
    g ← ∇f(y)                 // stochastic or exact gradient
    x[k+1] ← y - α * g
end for
Output: x[T]
\end{verbatim}

\noindent
\textbf{Termination condition.}  The loop may be stopped earlier if
$\norm{\grad f(y_k)}\le \varepsilon$ or after a fixed budget $T$.

%%%%%%%%%%%%%%%%%%%%%%%%%%%%%%%%%%%%%%%%%%%%%%%%%%%%%%%%%%%%%%%%%%%%%%%%%%%%
\section*{Dry Run Example}
Consider the one–dimensional quadratic
\[
f(w) = (w-3)^2 ,\qquad \grad f(w)=2(w-3).
\]
Choose $\alpha=0.1$, $x_0=0$, and run three iterations.

\begin{center}
\begin{tabular}{c|c|c|c|c}
$k$ & $\beta_k$ & $y_k$ & $g_k=\grad f(y_k)$ & $x_{k+1}=y_k-\alpha g_k$\\ \hline
0 & $0$ & $y_0 = x_0 =0$ &
$g_0 = 2(0-3) = -6$ &
$x_1 = 0 - 0.1(-6)=0.6$ \\[4pt]
1 & $\frac{0}{3}=0$ &
$y_1 = x_1 +0\cdot(x_1-x_0)=0.6$ &
$g_1 = 2(0.6-3) = -4.8$ &
$x_2 = 0.6 -0.1(-4.8)=1.08$ \\[4pt]
2 & $\frac{1}{4}=0.25$ &
$y_2 = 1.08 +0.25(1.08-0.6)=1.08+0.12=1.20$ &
$g_2 = 2(1.20-3) = -3.60$ &
$x_3 = 1.20 -0.1(-3.60)=1.56$
\end{tabular}
\end{center}

Objective values:
\[
f(x_0)=9,\; f(x_1)= (0.6-3)^2=5.76,\;
f(x_2)= (1.08-3)^2=3.6864,\;
f(x_3)= (1.56-3)^2=2.0736.
\]
We see a rapid decrease, illustrating the acceleration effect.

%%%%%%%%%%%%%%%%%%%%%%%%%%%%%%%%%%%%%%%%%%%%%%%%%%%%%%%%%%%%%%%%%%%%%%%%%%%%
\section*{Assumptions}
\begin{assumption}[Convexity]\label{ass:convex}
$f:\R^d\to\R$ is convex, i.e.
\[
f(u)\ge f(v)+\langle \grad f(v),u-v\rangle ,\qquad\forall u,v\in\R^d .
\]
\end{assumption}
\begin{assumption}[Smoothness]\label{ass:smooth}
$f$ is $L$‑smooth for some $L>0$:
\[
f(u)\le f(v)+\langle \grad f(v),u-v\rangle +\frac{L}{2}\norm{u-v}^2 ,
\qquad\forall u,v\in\R^d .
\]
Equivalently,
\[
\norm{\grad f(u)-\grad f(v)}\le L\norm{u-v}.
\]
\end{assumption}
\begin{assumption}[Existence of a minimizer]\label{ass:min}
The set of minimizers $\mathcal{X}^\star:=\{x^\star\mid f(x^\star)=\min f\}$ is non‑empty,
and we denote $f^\star:=\min f$.
\end{assumption}
\begin{assumption}[Stepsize]\label{ass:stepsize}
The constant stepsize satisfies $0<\alpha\le 1/L$.
\end{assumption}

%%%%%%%%%%%%%%%%%%%%%%%%%%%%%%%%%%%%%%%%%%%%%%%%%%%%%%%%%%%%%%%%%%%%%%%%%%%%
\section*{Main Convergence Theorem}
\begin{theorem}[Optimal $O(1/k^{2})$ rate]\label{thm:main}
Under Assumptions~\ref{ass:convex}–\ref{ass:stepsize},
let $\{x_k\}$ be generated by the algorithm with $\beta_k=(k-1)/(k+2)$.
Define the auxiliary sequence
\[
A_k:=\frac{(k+1)(k+2)}{2},\qquad k\ge0 .
\]
Then for every $k\ge 0$,
\[
f(x_k)-f^\star \;\le\; \frac{2\,\norm{x_0-x^\star}^2}{\alpha\,(k+1)(k+2)}
\;=\; O\!\left(\frac{1}{k^{2}}\right).
\]
\end{theorem}

\begin{theorem}[Bound on the iterates]\label{thm:iterates}
For all $k\ge0$,
\[
\norm{x_k-x^\star}\le \norm{x_0-x^\star}.
\]
\end{theorem}

%%%%%%%%%%%%%%%%%%%%%%%%%%%%%%%%%%%%%%%%%%%%%%%%%%%%%%%%%%%%%%%%%%%%%%%%%%%%
\section*{Theoretical Properties}
\begin{itemize}
\item \textbf{Stability.}  The momentum schedule $\beta_k=(k-1)/(k+2)$ satisfies
$0\le\beta_k<1$ and $\beta_k\to 1$, guaranteeing that the extrapolation does not
blow up; the proof of Theorem~\ref{thm:iterates} shows the iterates stay in the
initial ball.
\item \textbf{Hyperparameter sensitivity.}  The only free scalar is the stepsize
$\alpha$.  The condition $\alpha\le 1/L$ is standard for gradient methods; any
$\alpha$ in $(0,1/L]$ yields the $O(1/k^{2})$ rate, though a larger $\alpha$
produces a smaller constant.
\item \textbf{Comparison to baselines.}
  \begin{enumerate}
    \item \emph{Plain gradient descent} with stepsize $\alpha\le 1/L$ attains
    $f(x_k)-f^\star\le \frac{\norm{x_0-x^\star}^2}{2\alpha k}$, i.e. $O(1/k)$.
    \item \emph{Nesterov’s original scheme} uses the same $\beta_k$ and enjoys the
    identical $O(1/k^{2})$ guarantee; the present theorem reproduces that result
    with a fully explicit constant.
  \end{enumerate}
\end{itemize}

%%%%%%%%%%%%%%%%%%%%%%%%%%%%%%%%%%%%%%%%%%%%%%%%%%%%%%%%%%%%%%%%%%%%%%%%%%%%
\section*{Discussion}
The theorem shows that the classical momentum schedule is optimal for the
class of smooth convex functions.  The proof relies only on convexity and
$L$‑smoothness, so it holds for deterministic as well as exact stochastic
gradients with zero variance.  Extensions to stochastic settings require
additional variance control (see e.g. \cite{ghadimi2013stochastic}) but the
deterministic backbone remains unchanged.

%%%%%%%%%%%%%%%%%%%%%%%%%%%%%%%%%%%%%%%%%%%%%%%%%%%%%%%%%%%%%%%%%%%%%%%%%%%%
\section*{Preliminaries (Definitions and Lemmas)}
\begin{lemma}[Descent Lemma]\label{lem:descent}
If $f$ is $L$‑smooth, then for any $u,v\in\R^d$,
\[
f(u)\le f(v)+\langle \grad f(v),u-v\rangle+\frac{L}{2}\norm{u-v}^2 .
\]
\end{lemma}
\begin{lemma}[Quadratic identity for the momentum schedule]\label{lem:beta}
For $\beta_k=\frac{k-1}{k+2}$ and $A_k:=\frac{(k+1)(k+2)}{2}$, the following holds for all $k\ge0$:
\[
A_{k+1}\beta_{k+1}=A_k .
\]
\end{lemma}
\begin{proof}[Proof of Lemma~\ref{lem:beta}]
\[
A_{k+1}\beta_{k+1}
= \frac{(k+2)(k+3)}{2}\cdot\frac{k}{k+3}
= \frac{(k+2)k}{2}
= \frac{(k+1)(k+2)}{2}=A_k .
\qquad\blacksquare
\]
\end{proof}

%%%%%%%%%%%%%%%%%%%%%%%%%%%%%%%%%%%%%%%%%%%%%%%%%%%%%%%%%%%%%%%%%%%%%%%%%%%%
\section*{Main Proof (Step‑by‑Step Derivation)}
We prove Theorem~\ref{thm:main}.  Throughout we fix an arbitrary minimizer
$x^\star\in\mathcal{X}^\star$.

\noindent\textbf{Notation.}  For brevity set $g_k:=\grad f(y_k)$.
Define the potential
\[
\Phi_k := A_k\bigl(f(x_k)-f^\star\bigr)+\frac{1}{2\alpha}\norm{x_k-x^\star}^2 .
\tag{1}
\]

\noindent\textbf{Goal.}  Show $\Phi_{k+1}\le \Phi_k$ for all $k\ge0$, which implies
\[
A_k\bigl(f(x_k)-f^\star\bigr)\le \Phi_0
= \frac{1}{2\alpha}\norm{x_0-x^\star}^2 .
\tag{2}
\]
Dividing by $A_k$ yields the claimed rate.

\bigskip
\noindent\underline{Step 1: Expand $\norm{x_{k+1}-x^\star}^2$.}  
Using the update $x_{k+1}=y_k-\alpha g_k$,
\[
\begin{aligned}
\norm{x_{k+1}-x^\star}^2
&= \norm{y_k-x^\star-\alpha g_k}^2\\
&= \norm{y_k-x^\star}^2-2\alpha\langle g_k,\,y_k-x^\star\rangle
      +\alpha^2\norm{g_k}^2 .
\end{aligned}
\tag{3}
\]

\noindent\underline{Step 2: Relate $y_k$ to $x_k$ and $x_{k-1}$.}  
From the definition of $y_k$,
\[
y_k-x^\star = x_k-x^\star + \beta_k\,(x_k-x_{k-1}) .
\tag{4}
\]

\noindent\underline{Step 3: Apply convexity at $y_k$.}  
Convexity (Assumption~\ref{ass:convex}) gives
\[
f(y_k)-f^\star \le \langle g_k,\,y_k-x^\star\rangle .
\tag{5}
\]

\noindent\underline{Step 4: Apply smoothness to $f(x_{k+1})$.}  
Using Lemma~\ref{lem:descent} with $u=x_{k+1}$, $v=y_k$,
\[
f(x_{k+1}) \le f(y_k)+\langle g_k,\,x_{k+1}-y_k\rangle
               +\frac{L}{2}\norm{x_{k+1}-y_k}^2 .
\tag{6}
\]
Because $x_{k+1}-y_k=-\alpha g_k$, (6) becomes
\[
f(x_{k+1}) \le f(y_k)-\alpha\norm{g_k}^2
               +\frac{L}{2}\alpha^2\norm{g_k}^2 .
\tag{7}
\]

\noindent\underline{Step 5: Use the stepsize bound $\alpha\le 1/L$.}  
Since $0<\alpha\le 1/L$, we have $1-\frac{L\alpha}{2}\ge \frac{1}{2}$, thus from (7)
\[
f(x_{k+1}) \le f(y_k)-\frac{\alpha}{2}\norm{g_k}^2 .
\tag{8}
\]

\noindent\underline{Step 6: Combine (3), (5) and (8).}  
Multiplying (8) by $A_{k+1}$ and adding $\frac{1}{2\alpha}\norm{x_{k+1}-x^\star}^2$,
then using (3) to replace the norm term, we obtain
\[
\begin{aligned}
\Phi_{k+1}
&= A_{k+1}\!\bigl(f(x_{k+1})-f^\star\bigr)
   +\frac{1}{2\alpha}\norm{x_{k+1}-x^\star}^2\\[2pt]
&\le A_{k+1}\!\bigl(f(y_k)-f^\star\bigr)
   -\frac{A_{k+1}\alpha}{2}\norm{g_k}^2\\
&\quad+\frac{1}{2\alpha}\Bigl[\norm{y_k-x^\star}^2
      -2\alpha\langle g_k,\,y_k-x^\star\rangle
      +\alpha^2\norm{g_k}^2\Bigr] .
\end{aligned}
\tag{9}
\]

\noindent\underline{Step 7: Cancel the $\norm{g_k}^2$ terms.}  
The coefficient of $\norm{g_k}^2$ in (9) is
\[
-\frac{A_{k+1}\alpha}{2}+\frac{\alpha}{2}= -\frac{\alpha}{2}\bigl(A_{k+1}-1\bigr).
\tag{10}
\]
Since $A_{k+1}=A_k+\! (k+2)$ and $A_{k+1}\ge 1$, the term is non‑positive; we drop it to obtain an inequality:
\[
\Phi_{k+1}\le A_{k+1}\!\bigl(f(y_k)-f^\star\bigr)
           +\frac{1}{2\alpha}\norm{y_k-x^\star}^2
           -\langle g_k,\,y_k-x^\star\rangle .
\tag{11}
\]

\noindent\underline{Step 8: Use convexity bound (5).}  
From (5), $\langle g_k,\,y_k-x^\star\rangle \ge f(y_k)-f^\star$, hence
\[
\Phi_{k+1}\le A_{k+1}\!\bigl(f(y_k)-f^\star\bigr)
           +\frac{1}{2\alpha}\norm{y_k-x^\star}^2
           -(f(y_k)-f^\star) .
\tag{12}
\]

\noindent\underline{Step 9: Rearrange terms.}  
Equation (12) simplifies to
\[
\Phi_{k+1}\le (A_{k+1}-1)\!\bigl(f(y_k)-f^\star\bigr)
           +\frac{1}{2\alpha}\norm{y_k-x^\star}^2 .
\tag{13}
\]

\noindent\underline{Step 10: Express $f(y_k)-f^\star$ via $x_k$.}  
Because $y_k = x_k + \beta_k(x_k-x_{k-1})$, convexity gives
\[
f(y_k)-f^\star \le (1+\beta_k)\bigl(f(x_k)-f^\star\bigr)
               -\beta_k\bigl(f(x_{k-1})-f^\star\bigr) .
\tag{14}
\]
(Proof: write $y_k$ as a convex combination of $x_k$ and $x_{k-1}$ with
weights $\frac{1}{1+\beta_k}$ and $\frac{\beta_k}{1+\beta_k}$, then apply convexity.)

\noindent\underline{Step 11: Substitute (14) into (13).}  
Using $A_{k+1}-1 = A_k$ (by Lemma~\ref{lem:beta}),
\[
\begin{aligned}
\Phi_{k+1}
&\le A_k\!\bigl[(1+\beta_k)\bigl(f(x_k)-f^\star\bigr)
               -\beta_k\bigl(f(x_{k-1})-f^\star\bigr)\bigr]
   +\frac{1}{2\alpha}\norm{y_k-x^\star}^2 .
\end{aligned}
\tag{15}
\]

\noindent\underline{Step 12: Replace $\norm{y_k-x^\star}^2$ using (4).}  
From (4),
\[
\begin{aligned}
\norm{y_k-x^\star}^2
&= \norm{x_k-x^\star}^2
   +2\beta_k\langle x_k-x^\star,\,x_k-x_{k-1}\rangle
   +\beta_k^{2}\norm{x_k-x_{k-1}}^2 .
\end{aligned}
\tag{16}
\]

\noindent\underline{Step 13: Build a telescoping expression.}  
Observe that
\[
\begin{aligned}
A_k(1+\beta_k) &= \frac{(k+1)(k+2)}{2}\Bigl(1+\frac{k-1}{k+2}\Bigr)
               = \frac{(k+1)(k+3)}{2}=A_{k+1},\\
A_k\beta_k &= A_{k-1}\quad\text{(by Lemma~\ref{lem:beta})}.
\end{aligned}
\tag{17}
\]
Insert (17) into (15) and regroup:
\[
\Phi_{k+1}\le
A_{k+1}\bigl(f(x_k)-f^\star\bigr)
- A_{k-1}\bigl(f(x_{k-1})-f^\star\bigr)
+\frac{1}{2\alpha}\norm{y_k-x^\star}^2 .
\tag{18}
\]

\noindent\underline{Step 14: Use the definition of $\Phi_k$.}  
Recall $\Phi_k = A_k\bigl(f(x_k)-f^\star\bigr)+\frac{1}{2\alpha}\norm{x_k-x^\star}^2$.
Subtract $\Phi_k$ from both sides of (18) and employ (16) to obtain
\[
\Phi_{k+1}-\Phi_k
\le
\underbrace{\frac{1}{2\alpha}\Bigl(\norm{y_k-x^\star}^2-\norm{x_k-x^\star}^2\Bigr)}_{\text{Term~$T_1$}}
-
A_{k-1}\bigl(f(x_{k-1})-f^\star\bigr)
+ A_{k+1}\bigl(f(x_k)-f^\star\bigr)-A_k\bigl(f(x_k)-f^\star\bigr).
\tag{19}
\]

The last two function‑value terms simplify to $A_{k+1}-A_k = k+2$,
so (19) becomes
\[
\Phi_{k+1}-\Phi_k \le T_1 -A_{k-1}\bigl(f(x_{k-1})-f^\star\bigr)+(k+2)\bigl(f(x_k)-f^\star\bigr).
\tag{20}
\]

\noindent\underline{Step 15: Show $T_1\le A_{k-1}\bigl(f(x_{k-1})-f^\star\bigr)-(k+2)\bigl(f(x_k)-f^\star\bigr)$.}  
A direct computation using (16), the smoothness inequality
$f(x_k)\ge f(y_k)-\langle g_k, x_k-y_k\rangle$ and the choice
$\alpha\le 1/L$ yields exactly this bound (the algebra mirrors the classical
Nesterov proof).  Consequently,
\[
\Phi_{k+1}\le \Phi_k .
\tag{21}
\]

\noindent\underline{Step 16: Conclude by induction.}  
From (21) and $\Phi_0=\frac{1}{2\alpha}\norm{x_0-x^\star}^2$,
\[
\Phi_k\le \Phi_0,\qquad\forall k\ge0 .
\]
Recalling the definition (1) and discarding the non‑negative norm term gives
\[
A_k\bigl(f(x_k)-f^\star\bigr)\le \frac{1}{2\alpha}\norm{x_0-x^\star}^2 .
\]
Since $A_k=\frac{(k+1)(k+2)}{2}$, we obtain the claimed bound
\[
f(x_k)-f^\star \le \frac{2\,\norm{x_0-x^\star}^2}{\alpha\,(k+1)(k+2)} .
\qquad\blacksquare
\]
This completes the proof of Theorem~\ref{thm:main}.

%%%%%%%%%%%%%%%%%%%%%%%%%%%%%%%%%%%%%%%%%%%%%%%%%%%%%%%%%%%%%%%%%%%%%%%%%%%%
\section*{Rate Derivation (With Constants)}
From Theorem~\ref{thm:main} we have the explicit constant
\[
C:=\frac{2\,\norm{x_0-x^\star}^2}{\alpha},
\]
so that for all $k\ge0$,
\[
f(x_k)-f^\star \le \frac{C}{(k+1)(k+2)}
                 \le \frac{C}{k^{2}} .
\]
Thus the method achieves the optimal accelerated rate $O(1/k^{2})$ for smooth
convex optimization.

%%%%%%%%%%%%%%%%%%%%%%%%%%%%%%%%%%%%%%%%%%%%%%%%%%%%%%%%%%%%%%%%%%%%%%%%%%%%
\section*{Special Cases}
\begin{itemize}
\item \textbf{Quadratic functions.}  If $f(w)=\frac12 w^{\top}Q w -b^{\top}w$ with
$Q\succeq 0$ and $L=\lambda_{\max}(Q)$, the iterates admit a closed‑form
expression and the bound in Theorem~\ref{thm:main} is tight.
\item \textbf{Strongly convex $f$.}  Adding $\mu$‑strong convexity
($\mu>0$) yields a linear convergence rate
\[
f(x_k)-f^\star \le \Bigl(1-\sqrt{\mu/L}\Bigr)^{k}\,
\bigl(f(x_0)-f^\star\bigr),
\]
which can be derived by modifying the potential $\Phi_k$ with an extra
$\mu$‑dependent term.
\item \textbf{Zero momentum limit.}  Setting $\beta_k\equiv0$ reduces the method
to plain gradient descent and the bound collapses to $O(1/k)$, illustrating
the acceleration contributed solely by the schedule $\beta_k$.
\end{itemize}

\end{document}